\chap{热处理温度场物理模型}

本文研究给定初态、炉温动作及边界材料参数下的传热过程。数值求解器生成的温度轨迹构成离散状态转移的学习数据，传热方程、物性函数和边界条件同时约束温度预测、边界辨识与闭环控制的适用范围。

\sect{研究对象与基本假设}

主要研究对象为两侧受热条件相同的 C45/AISI 1045 类中碳钢平板。$L$ 表示平板总厚度，完整一维截面为 $x\in[0,L]$，中心位于 $x=L/2$。均匀初温与相同的两侧边界使温度场关于中心面对称。模型假设材料宏观各向同性，截面内无体热源，炉内对流和辐射沿表面均匀分布。研究温区覆盖加热和随后冷却过程。导热系数与比热随温度变化，密度取 \SI{7870}{\kilogram\per\cubic\metre}。

三维研究对象为固定长方体试块，其六个表面采用均匀的对流--辐射条件，材料物性模型与一维平板相同。一维环境承担温度推演、边界辨识与闭环控制研究，三维实验考察固定几何上的温度场推演。两者的状态方程均不包含相变动力学、组织演化、热应力与变形。温变比热表反映部分高温热容变化，但不单独计算潜热，温度预测结果也不直接给出组织或力学性能。

\sect{控制方程}

无内热源的一维非稳态导热方程为
\begin{equation}
  \rho c_p(T)\frac{\partial T}{\partial t}
  =\frac{\partial}{\partial x}\left[k(T)\frac{\partial T}{\partial x}\right],
  \qquad 0<x<L.
  \label{eq:heat_pde}
\end{equation}
式中，$T$ 为温度，$\rho$ 为密度，$c_p(T)$ 为比热容，$k(T)$ 为导热系数。

两侧表面采用对流与灰体辐射联合边界。以指向工件内部的热流为正，有
\begin{equation}
  q_s=h(\Tinfty-T_s)
  +\varepsilon\sigma\left(T_{\infty,K}^{4}-T_{s,K}^{4}\right),
  \label{eq:mixed_boundary}
\end{equation}
其中 $T_s$ 为表面温度，$\Tinfty$ 为炉气与等效炉壁温度，$h$ 为对流换热系数，$\varepsilon$ 为表面发射率，$\sigma$ 为 Stefan--Boltzmann 常数，下标 K 表示开尔文温度。

四次方差可因式分解为
\begin{equation}
  T_{\infty,K}^{4}-T_{s,K}^{4}
  =(T_{\infty,K}-T_{s,K})(T_{\infty,K}+T_{s,K})
  (T_{\infty,K}^{2}+T_{s,K}^{2}).
\end{equation}
定义
\begin{equation}
  h_{\mathrm{rad}}(T_s,\Tinfty)
  =\sigma(T_{\infty,K}+T_{s,K})(T_{\infty,K}^{2}+T_{s,K}^{2}),
  \label{eq:hrad}
\end{equation}
边界热流可以写为
\begin{equation}
  q_s=\left[h+\varepsilon h_{\mathrm{rad}}(T_s,\Tinfty)\right](\Tinfty-T_s).
  \label{eq:heff_flux}
\end{equation}
式\eqref{eq:heff_flux}中的边界热流由等效换热系数与温差的乘积确定。等效系数随表面温度和炉温变化，是边界可辨识性分析中的状态、动作相关参数组合。

\sect{温变物性与参数空间}

高温物性采用公开 AISI 1045 数据表的分段线性插值\cite{sanchez2023splat}。Carollo 等的低温实验反演结果用于核对低温段变化趋势\cite{carollo2019properties}。表\ref{tab:properties}列出计算所用节点值。原始数据在 \SI{600}{\degreeCelsius} 至 \SI{800}{\degreeCelsius} 间给出明显的比热变化，本文按表值插值，并将其作为经验等效热容曲线。相变附近的表观热容效应可能包含在经验数据中，但数据源未将其与显热、潜热分别标定。当前模型不建立显式相变动力学，也不据此解释组织转变。

\begin{table}[H]
  \centering
  \caption{C45/AISI 1045 类中碳钢温变物性}
  \label{tab:properties}
  \begin{tabular}{c*{6}{r}}
    \toprule
    温度/\si{\degreeCelsius} & 0 & 200 & 400 & 600 & 800 & 1000 \\
    \midrule
    $c_p$/\si{\joule\per\kilogram\per\kelvin} & 510 & 600 & 725 & 900 & 562 & 600 \\
    $k$/\si{\watt\per\metre\per\kelvin} & 45 & 41 & 36 & 32 & 27 & 23 \\
    \bottomrule
  \end{tabular}
\end{table}

数据生成时对 $k(T)$ 和 $c_p(T)$ 分别乘以尺度因子 $s_k$ 与 $s_c$，以考察模型对物性变化的响应。两个尺度的训练范围均为 0.95 至 1.05。平板厚度在 \SIrange{10}{40}{\milli\metre} 内采样，对流换热系数在 \SIrange{10}{60}{\watt\per\square\metre\per\kelvin} 内采样，发射率范围为 $0.65$ 至 $0.90$。初始温度为 \SIrange{15}{40}{\degreeCelsius}，炉温上限为 \SI{950}{\degreeCelsius}。这些范围界定了学习与评价所用的受控数值环境，不代表特定炉次的标定值。

在 $L=\SI{20}{\milli\metre}$、$h=\SI{25}{\watt\per\square\metre\per\kelvin}$、$\varepsilon=0.75$、$\Tinfty=\SI{850}{\degreeCelsius}$ 的算例中，辐射等效系数由初始约 \SI{81.62}{\watt\per\square\metre\per\kelvin} 增至 \SI{146.85}{\watt\per\square\metre\per\kelvin}。辐射贡献高于对流项。300 s 后，联合边界给出的中心温度为 \SI{475.59}{\degreeCelsius}，纯对流计算值为 \SI{154.99}{\degreeCelsius}。本文的高温工况保留完整的对流--辐射边界。

\sect{数值离散与参考求解器}
\label{sec:physical-discretization}

空间离散采用基于节点控制体能量平衡的规则网格守恒通量差分，时间推进分别采用冻结系数的半隐式格式和自适应积分。炉温动作 $a_n$ 表示区间 $[t_n,t_{n+1})$ 内保持不变的设定值，数值求解、网络输入和物理残差均使用这一动作定义。

一维训练数据生成器采用 41 个等距节点，时间步长为 \SI{1}{\second}。界面导热系数取相邻节点旧时刻导热系数的算术平均。内部节点 $i$ 的离散平衡写为
\begin{equation}
  \rho c_{p,i}^{n}\frac{T_i^{n+1}-T_i^n}{\Delta t}
  =\frac{1}{\Delta x}
  \left[
    k_{i+1/2}^{n}\frac{T_{i+1}^{n+1}-T_i^{n+1}}{\Delta x}
    -k_{i-1/2}^{n}\frac{T_i^{n+1}-T_{i-1}^{n+1}}{\Delta x}
  \right].
  \label{eq:discrete_balance}
\end{equation}
物性和辐射等效系数由旧状态计算，新状态通过隐式线性系统求解。以左侧半控制体为例，表面离散式为
\begin{equation}
  T_1^{n+1}-T_1^n
  =\frac{2\Delta t}{\rho c_{p,1}^{n}\Delta x^2}
  \left[
    k_{3/2}^{n}(T_2^{n+1}-T_1^{n+1})
    +\Delta x\Heff_1^n(\Tinfty^n-T_1^{n+1})
  \right],
  \label{eq:surface_discrete_balance}
\end{equation}
其中 $\Heff_1^n=h^n+\varepsilon^nh_{\mathrm{rad}}(T_1^n,\Tinfty^n)$，右侧表面采用镜像形式。系数中的表面温度取旧状态，换热温差中的表面温度取新状态。一维数据生成与训练物理残差共用这一半隐式能量平衡格式。

一维交叉求解器验证与控制环境使用 81 节点空间网格。方法线离散所得的常微分方程组由 SciPy 自适应后向差分公式（backward differentiation formula，BDF）积分器推进。参考求解器与训练生成器共用连续方程和物性函数，采用不同的空间分辨率及时间积分方法，并在积分过程中更新温度相关的物性和辐射热流。参考求解器采用相对容差 $10^{-6}$、绝对温度容差 $10^{-7}$~\si{\degreeCelsius}、最大内部步长 \SI{0.25}{\second}。161 节点与更小时间步的收敛检查分别得到平均 RMSE \SI{0.00248}{\degreeCelsius} 和 \SI{0.00425}{\degreeCelsius}。训练生成器与 BDF 参考轨迹的整体 RMSE 为 \SI{0.166}{\degreeCelsius}。这些一维 BDF 参考轨迹不用于神经模型训练，提供数值求解方式改变后的评价依据。

三维实验在规则长方体网格上沿三个空间方向计算守恒通量，基础网格 $9\times7\times5$ 与加密网格 $17\times13\times9$ 分别包含 315 和 1989 个温度状态。三维轨迹直接由方法线与 BDF 生成，并按工况划分后用于训练、验证和测试。三维训练采用半隐式离散残差，与 BDF 目标轨迹存在时间离散差异。数值求解结果的解释限于既定传热模型，交叉求解器一致性和网格检查不能替代真实炉体实验。

\sect{受控数值数据集}
\label{sec:dataset-design}

数据集以一组初态、炉温序列和边界材料参数定义一条物理工况，每条轨迹包含 300 个状态转移。训练、验证和测试按完整轨迹划分，避免相邻时刻进入不同数据划分。训练集用于更新网络及确定归一化尺度，验证集用于选择模型，测试集用于比较留出工况上的误差。

一维训练及分布内（ID）测试包含线性升温保温和两段加热。分布外（OOD）评价分别改变控制曲线、参数取值或边界时间结构，以区分误差来源。控制 OOD 使用炉温阶跃及升温后冷却；参数 OOD 保持控制曲线类型不变，将 $h$、$\varepsilon$、$s_k$ 和 $s_c$ 推到表\ref{tab:ood_ranges}所列区间。每个单因素类别包含低侧和高侧各 10 条轨迹，组合类别对四个参数分别随机选择低侧或高侧区间。

\begin{table}[H]
  \centering
  \caption{训练参数范围与参数 OOD 测试范围}
  \label{tab:ood_ranges}
  \begin{tabular}{lccc}
    \toprule
    参数 & 训练范围 & OOD 低侧 & OOD 高侧 \\
    \midrule
    $h$/\si{\watt\per\square\metre\per\kelvin} & $10$--$60$ & $5$--$9$ & $65$--$80$ \\
    $\varepsilon$ & $0.65$--$0.90$ & $0.45$--$0.60$ & $0.92$--$0.98$ \\
    $s_k$ & $0.95$--$1.05$ & $0.80$--$0.90$ & $1.10$--$1.20$ \\
    $s_c$ & $0.95$--$1.05$ & $0.80$--$0.90$ & $1.10$--$1.20$ \\
    \bottomrule
  \end{tabular}
\end{table}

动态边界 OOD 使 $h(t)$ 与 $\varepsilon(t)$ 在单条轨迹内变化，同时保持瞬时值位于训练范围。$h$ 的测试形式包括余弦平滑过渡、在轨迹 35\%--65\% 区间内随机选取时刻发生的阶跃和正弦变化；$\varepsilon$ 采用指数增长或阶跃变化。联合类别组合正弦 $h(t)$ 与增长型 $\varepsilon(t)$。训练数据中的边界参数在每条轨迹内保持常数，所以该测试隔离了边界时间结构的变化。

数值参考与学习模型共享传热方程、材料函数和几何假设，二者的误差比较描述学习模型对指定热动力学环境的逼近程度。边界与内部状态的估计误差及其控制影响由测量驱动的闭环实验评价。
